\chapter{Vision of the solution}
\label{cha:vision}

This chapter states the requirements the design in \Cref{cha:design} must satisfy, before any library or model is chosen. Two shape the architecture more than the rest: the system may never report a nutrient figure it cannot trace to a database entry, and the grounding and review mechanisms under study must be independently switchable so that a measured difference can be attributed to one toggled mechanism. The remaining requirements fix the one-day session shape, the recipe-and-grams contract, and the latency bound that later rules out an unbounded review loop.

\section{Usage scenario}
\label{sec:vision-usage-scenario}

A user's profile -- age, sex, body measurements, activity level, a dietary goal, any allergies or diagnosed conditions, and any standing food preferences or dislikes -- is already on file. At the start of a session the system asks the user to describe, in free text, anything specific they want for that day (an ingredient to use up, a craving, a time constraint) or to proceed with no additional request. From this input the system produces one complete plan for one day: a full day's meals, each with a name, a step-by-step preparation method, and the exact grams of every ingredient. The system does not conduct a multi-turn conversation that edits a previous plan, nor does it remember what the user ate afterwards; each session produces one plan from a stated set of inputs. Both omissions concern carrying state between sessions and introduce no mechanism this thesis measures, so both are left to future work (\Cref{sec:conclusions-future-work}).

\section{Requirements}
\label{sec:vision-requirements}

The requirements are grouped by the kind of claim each makes: what a plan must always be, what it must try to be, and what the system must permit in order to be studied.

\begin{description}
    \freq{req:one-day-plan} The system shall accept a fixed user profile (demographic data, dietary goal, allergies, medical conditions, food preferences) and an optional free-form session request, and shall produce a single, complete plan covering every meal of one day.
    \freq{req:recipes-and-grams} Every meal in the plan shall include a named recipe with preparation instructions detailed enough for a home cook to follow, and a list of ingredients with their exact quantities in grams.
    \freq{req:no-invented-nutrients} Every ingredient quantity in the plan shall correspond to a real entry in a food knowledge base, not to a value invented by the system; the system shall never report a nutrient figure that cannot be traced to a knowledge base entry and a specific quantity.
\end{description}

Two mechanisms enforce these before delivery. The output schema rejects a plan with no meals, an unnamed recipe, an empty preparation method or a meal without portions, which is what \ref{req:one-day-plan}~and \ref{req:recipes-and-grams}~reduce to once the plan is a typed structure. A separate validator rejects a plan naming a food that fails to resolve in the database, which is \ref{req:no-invented-nutrients}. That last guarantee is absolute, not statistical, and it covers every nutrient figure the system reports as data; a number the model writes into its free-text explanation is never parsed as nutrition data and is outside the guarantee (\Cref{sec:design-no-invented-foods}).

The next two are goals, not guarantees, because both rest on the model's judgement about food. An exclusion under \ref{req:constraints} must be read off a food's name, category and ingredient text; the data carry no reliable allergen or diet-pattern tag to filter on. A target under \ref{req:macro-targets} can only be aimed at. Both are measured per plan instead of being enforced: the hard half of \ref{req:constraints} by the safety score of \Cref{def:safety-score}, its soft half by the soft score of \Cref{def:soft-score}, and \ref{req:macro-targets} by the accuracy results of \Cref{cha:results}.

\begin{description}
    \freq{req:constraints} No meal in the plan shall include a food that conflicts with a hard exclusion stated in the user's profile (an allergen, or a diet pattern such as vegetarianism), and the plan shall take reasonable account of soft preferences (liked and disliked foods, a session request) where doing so does not conflict with a hard exclusion or a medical constraint.
    \freq{req:macro-targets} The plan shall target the energy and macronutrient levels implied by the user's profile and dietary goal; the targets themselves, and the delivered totals they are compared against, shall be computed independently of the language model's own claims about either.
\end{description}

The last three constrain the system, not the plan: two make the ablation study possible, and the third caps how long a session may take.

\begin{description}
    \nfreq{req:ablatable} The system's behaviour shall be reproducible and measurable well enough to support a controlled ablation study: individual grounding and review mechanisms shall be independently switchable without altering anything else about the pipeline, so that a difference in outcome between two configurations can be attributed to the mechanism that was toggled.
    \nfreq{req:provider-agnostic} The system shall be usable with any of several LLM providers, and shall expose the reasoning-effort setting a provider offers as a configurable parameter, so that the effect of model capability on plan quality can be measured along both axes without being tied to a single vendor's model family.
    \nfreq{req:interactive-runtime} Every user-visible run of the system shall complete within a time compatible with an interactive session, since a plan that takes many minutes to produce would not be a viable substitute for existing planning workflows.
\end{description}

The system succeeds if it satisfies these requirements -- the enforceable ones holding on every run, the rest measured and reported -- and if the totaller and the critique-and-refine loop are switchable in isolation so that every hypothesis of \Cref{sec:introduction-goal} gets a verdict. A mechanism found to contribute nothing satisfies this condition as well as one found to help.

\section{Why food lookup is a fixed premise}
\label{sec:vision-food-lookup-premise}

An early design question was whether the ablation study should also test a variant with no food knowledge base access -- a language model reasoning about nutrient values purely from its training data. It was excluded, but for a narrower reason than it might appear. The quantitative metric this thesis reports, defined later in \Cref{def:macro-error}, is computed by the deterministic totaller from portions resolved against the real knowledge base, so a variant with no such access supplies nothing that metric can be computed from: its plan names foods without identifying them, and an unidentified food has no known composition to sum. Matching those names to entries after the fact would measure how well the post-hoc matcher guessed what the model meant, not the plan. That leaves only the model's own statement of its totals, which is not a measurement.

This rules out one of the two evaluation dimensions, not both. The judged dimension of \Cref{sec:methodology-qualitative} reads recipes and portions, so it could have scored an ungrounded plan perfectly well. What settled the question was the grid budget: an arm scored on one dimension and blank on the other could not be set beside the rest of \Cref{cha:results}, and a fifth variant would have added fifteen runs to a grid already at the limit of what this project could pay for. The ablation study in \Cref{cha:methodology} therefore treats the food knowledge base as always present, and varies only the deterministic totaller and the critique-and-refine loop; \Cref{sec:conclusions-future-work} notes that an ungrounded arm scored on the judged dimension alone remains open. \Cref{fig:diagram-variants} shows the four configurations this leaves.

\begin{figure}[htbp]
    \centering
    \begin{tikzpicture}[
    node distance=4mm,
    box/.style={draw, rounded corners, align=center, text width=42mm, minimum height=8mm, inner sep=1.5mm, font=\small},
    off/.style={box, dashed, draw=black!35, text=black!35},
    frame/.style={draw, rounded corners, inner sep=3mm},
    title/.style={font=\small\bfseries, align=center},
]

% --- baseline ---
\node[box] (b_agents) {Nutrition agent};
\node[box, below=of b_agents] (b_lookup) {Food knowledge database};
\node[off, below=of b_lookup] (b_tot) {Deterministic totaller};
\node[off, below=of b_tot] (b_ref) {Critique-and-refine loop};
\node[frame, fit=(b_agents)(b_ref)] (b_frame) {};
\node[title, above=1.5mm of b_frame.north] {baseline};

% --- totaller-only ---
\node[box, right=16mm of b_agents] (t_agents) {Nutrition agent};
\node[box, below=of t_agents] (t_lookup) {Food knowledge database};
\node[box, below=of t_lookup] (t_tot) {Deterministic totaller};
\node[off, below=of t_tot] (t_ref) {Critique-and-refine loop};
\node[frame, fit=(t_agents)(t_ref)] (t_frame) {};
\node[title, above=1.5mm of t_frame.north] {totaller-only};

% --- reflection-only ---
\node[box, below=18mm of b_ref] (r_agents) {Nutrition agent};
\node[box, below=of r_agents] (r_lookup) {Food knowledge database};
\node[off, below=of r_lookup] (r_tot) {Deterministic totaller};
\node[box, below=of r_tot] (r_ref) {Critique-and-refine loop};
\node[frame, fit=(r_agents)(r_ref)] (r_frame) {};
\node[title, above=1.5mm of r_frame.north] {reflection-only};

% --- full ---
\node[box, right=16mm of r_agents] (f_agents) {Nutrition agent};
\node[box, below=of f_agents] (f_lookup) {Food knowledge database};
\node[box, below=of f_lookup] (f_tot) {Deterministic totaller};
\node[box, below=of f_tot] (f_ref) {Critique-and-refine loop};
\node[frame, fit=(f_agents)(f_ref)] (f_frame) {};
\node[title, above=1.5mm of f_frame.north] {full};

\end{tikzpicture}

    \caption{The four-way ablation grid.}
    \label{fig:diagram-variants}
\end{figure}
